% Reusable English paper template. Save this file as UTF-8.
% Compile with pdfLaTeX, or run: latexmk -pdf arxiv_paper_template.tex
% Compile at least twice to update references and appendix contents.
%
% 使用方法：复制本文件，修改题目、作者和正文即可。
% 配色集中在下方“COLORS AND LINKS”部分。
% 文献采用作者—年份引用，如 Robbins and Monro (1951)，不用数字编号引用。
% 本文件自带示例参考文献和图框，不需要额外的 .bib、.bst 或图片文件。
% References 之后自动生成只含附录及小节的目录。

\documentclass[10pt,a4paper]{article}
\usepackage[a4paper,margin=1in]{geometry}
\usepackage[T1]{fontenc}
\usepackage{amsmath,amssymb,amsthm,mathtools,bm,mathrsfs}
\usepackage{graphicx,booktabs,array,microtype,aliascnt}
\usepackage{algorithm,algpseudocode}
\usepackage{xurl}
\usepackage[authoryear,round]{natbib}
\usepackage{xcolor}

% ================================================================
% COLORS AND LINKS -- 与 Muon 主稿保持一致，只需在这里修改颜色。
% ================================================================
% 文献引用和网址：深蓝色。
\definecolor{ReferenceBlue}{rgb}{0,0.19,0.56}
% 公式、定理、图表、算法、章节及附录引用：暗红色。
\definecolor{CrossReferenceRed}{rgb}{0.7,0,0}

\usepackage[
  colorlinks=true,
  linkcolor=CrossReferenceRed,
  citecolor=ReferenceBlue,
  urlcolor=ReferenceBlue,
  linktoc=all,
  pdfusetitle
]{hyperref}
\usepackage{bookmark}

% ================================================================
% THEOREM ENVIRONMENTS
% aliascnt 使共享编号的 lemma/corollary 等被 cleveref 正确识别。
% ================================================================
\newtheorem{theorem}{Theorem}[section]
\newaliascnt{lemma}{theorem}
\newtheorem{lemma}[lemma]{Lemma}
\aliascntresetthe{lemma}
\newaliascnt{proposition}{theorem}
\newtheorem{proposition}[proposition]{Proposition}
\aliascntresetthe{proposition}
\newaliascnt{corollary}{theorem}
\newtheorem{corollary}[corollary]{Corollary}
\aliascntresetthe{corollary}

\theoremstyle{definition}
\newtheorem{assumption}{Assumption}[section]
\newtheorem{definition}{Definition}[section]
\theoremstyle{remark}
\newtheorem{remark}{Remark}[section]

% cleveref 应放在 hyperref 和上述定理环境之后。
% 不使用 nameinlink：例如 Theorem 4.1 中仅编号标红，文字保持黑色。
\usepackage[noabbrev]{cleveref}
\crefname{algorithm}{Algorithm}{Algorithms}
\Crefname{algorithm}{Algorithm}{Algorithms}
\crefname{theorem}{Theorem}{Theorems}
\Crefname{theorem}{Theorem}{Theorems}
\crefname{lemma}{Lemma}{Lemmas}
\Crefname{lemma}{Lemma}{Lemmas}
\crefname{proposition}{Proposition}{Propositions}
\Crefname{proposition}{Proposition}{Propositions}
\crefname{corollary}{Corollary}{Corollaries}
\Crefname{corollary}{Corollary}{Corollaries}
\crefname{assumption}{Assumption}{Assumptions}
\Crefname{assumption}{Assumption}{Assumptions}
\crefname{definition}{Definition}{Definitions}
\Crefname{definition}{Definition}{Definitions}
\crefname{remark}{Remark}{Remarks}
\Crefname{remark}{Remark}{Remarks}
\crefname{section}{Section}{Sections}
\Crefname{section}{Section}{Sections}
\crefname{figure}{Figure}{Figures}
\Crefname{figure}{Figure}{Figures}
\crefname{table}{Table}{Tables}
\Crefname{table}{Table}{Tables}

\newcommand{\R}{\mathbb R}
\newcommand{\E}{\mathbb E}
\setlength{\emergencystretch}{2em}
\pagestyle{plain}
\raggedbottom

% ================================================================
% TITLE AND AUTHORS -- PDF metadata follows these automatically.
% ================================================================
\title{Paper Title}
\author{}
% 需要作者信息时，将上一行替换为例如：
% \author{Author Name\\Department, University\\
%   \texttt{author@example.com}}
\date{}

\begin{document}
% 保留此行：附录目录不显示正文各节。
\addtocontents{toc}{\protect\setcounter{tocdepth}{-1}}
\maketitle

\begin{abstract}
Replace this paragraph with the abstract of your paper. State the
problem, the assumptions, the main result and its scope. The examples
below demonstrate the reference colors and can be replaced directly.
\end{abstract}

\section{Introduction}

Introduce the problem and explain the main contribution. A literature
citation such as \citep{robbins1951} uses author--year labels in deep blue. Internal references
such as \cref{ass:smoothness,thm:descent}, \eqref{eq:descent}, and
\cref{alg:example,fig:example,tab:colors} have dark-red numbers.
The text of the paper, displayed equations and theorem headings
remain black.

\section{Related works}

Discuss the relevant literature and distinguish the assumptions and
conclusions. For a citation used as part of a sentence, write, for
example, \citet{robbins1951}.

\section{Problem setup}

Specify the objective, the domain and the algorithm. The Euclidean
norm is denoted by $\|\cdot\|$. The following assumption and algorithm
are illustrative examples rather than requirements of this template.

\begin{assumption}[$L$-smoothness]
\label{ass:smoothness}
The function $f:\R^d\to\R$ is differentiable, and there is $L>0$ such that
$\|\nabla f(x)-\nabla f(y)\|\le L\|x-y\|$ for all $x,y\in\R^d$.
\end{assumption}

\begin{algorithm}[htbp]
\caption{Gradient descent: an illustrative algorithm}
\label{alg:example}
\begin{algorithmic}[1]
\Require Initial point $x_0\in\R^d$, stepsize $\eta>0$, iterations $T\ge1$
\For{$t=0,\ldots,T-1$}
  \State $x_{t+1}\gets x_t-\eta\nabla f(x_t)$
\EndFor
\State \Return $x_T$
\end{algorithmic}
\end{algorithm}

\section{Main results}

\begin{theorem}[Illustrative descent result]
\label{thm:descent}
Under \cref{ass:smoothness}, if $0<\eta\le1/L$, the iterates of
\cref{alg:example} satisfy, for every $0\le t<T$,
\begin{equation}
 f(x_{t+1})\le f(x_t)-\frac{\eta}{2}\|\nabla f(x_t)\|^2.
 \label{eq:descent}
\end{equation}
\end{theorem}

The proof is in Appendix~\ref{app:proofs}. Use \verb|\cref{...}| for a
named reference and \verb|\eqref{...}| for an equation number in
parentheses. Use \verb|\citep{...}| or \verb|\citet{...}| for literature.

\section{Experiments}

Describe the experimental setup and report the findings. The examples
in \cref{fig:example,tab:colors} demonstrate captions and references.

\begin{figure}[htbp]
\centering
\fbox{\parbox[c][20mm][c]{0.75\linewidth}{\centering
  Figure placeholder\\[3pt]
  Replace this box with \texttt{\string\includegraphics}.}}
\caption{A self-contained figure placeholder. Figure captions remain black.}
\label{fig:example}
\end{figure}

\begin{table}[htbp]
\centering
\caption{The reference palette used throughout the template.}
\label{tab:colors}
\begin{tabular}{ll}
\toprule
Element & Color\\
\midrule
Internal cross-references & \textcolor{CrossReferenceRed}{Dark red}\\
Literature citations and URLs & \textcolor{ReferenceBlue}{Deep blue}\\
Body text, equations and captions & Black\\
\bottomrule
\end{tabular}
\end{table}

\section{Conclusion}

Summarize the results and their limitations. Replace the illustrative
content with your own paper while retaining the preamble and appendix
contents commands as needed.

% ================================================================
% REFERENCES -- single-file example, before the appendix contents.
% ================================================================
\clearpage
\begin{thebibliography}{9}
\bibitem[Robbins and Monro(1951)]{robbins1951}
H.~Robbins and S.~Monro.
\newblock A stochastic approximation method.
\newblock \emph{The Annals of Mathematical Statistics},
22(3):400--407, 1951.
\newblock \url{https://doi.org/10.1214/aoms/1177729586}.
\end{thebibliography}

% 如需 BibTeX，用下面两行替换整个 thebibliography 环境，并提供 .bib：
% plainnat 支持作者—年份引用；不要改成仅提供数字引用的参考文献样式。
% \bibliographystyle{plainnat}
% \bibliography{references}

% ================================================================
% APPENDICES -- printed contents and PDF sidebar entries in dark red.
% ================================================================
\clearpage
\appendix
\bookmarksetup{color=CrossReferenceRed}
\addtocontents{toc}{\protect\setcounter{tocdepth}{2}}
\begingroup
\color{CrossReferenceRed}
\hypersetup{linktoc=all}
\renewcommand{\contentsname}{Contents of the Appendices}
\phantomsection
\pdfbookmark[1]{Contents of the Appendices}{appendix-contents}
\tableofcontents
\endgroup
\clearpage
\numberwithin{equation}{section}

\section{Proofs}
\label{app:proofs}

\subsection{Proof of the illustrative theorem}

\begin{proof}[Proof of \cref{thm:descent}]
The smoothness assumption gives
\begin{equation}
 f(y)\le f(x)+\langle\nabla f(x),y-x\rangle
             +\frac{L}{2}\|y-x\|^2.
 \label{eq:appendix-smoothness}
\end{equation}
Substituting $x=x_t$ and $y=x_t-\eta\nabla f(x_t)$ into
\eqref{eq:appendix-smoothness}, we obtain
\[
 f(x_{t+1})\le f(x_t)
 -\eta\left(1-\frac{L\eta}{2}\right)\|\nabla f(x_t)\|^2.
\]
Since $L\eta\le1$, this proves \eqref{eq:descent}.
\end{proof}

\section{Additional details}

\subsection{Supplementary material}

Place further proofs, implementation details or additional results here.
The appendix contents update automatically as sections are added or
removed. Compile twice, or use \texttt{latexmk}, to refresh page numbers.

\end{document}
